% ============================================================
%  cap_potenze.tex — Le potenze
%  Matemagica — Scuole medie Canton Ticino (DECS)
% ============================================================

\chapter{Le potenze}

% --- Obiettivi di apprendimento --------------------------------
\section*{Obiettivi di apprendimento}
\begin{itemize}[leftmargin=1.5em]
  \item Comprendere il significato di \textbf{potenza} come moltiplicazione ripetuta.
  \item Conoscere e applicare le \textbf{proprietà delle potenze}.
  \item Riconoscere la precedenza dell'elevamento a potenza nelle espressioni aritmetiche.
  \item Leggere e scrivere numeri grandi utilizzando le \textbf{potenze di 10}.
  \item Comprendere i principi di base della \textbf{notazione scientifica}.
\end{itemize}

% ============================================================
\section{Dalla moltiplicazione alla potenza}
% ============================================================

\subsection{La riproduzione batterica}

I batteri sono microrganismi unicellulari che si riproducono dividendosi a metà:
da ogni divisione derivano due batteri, ognuno dei quali si divide di nuovo, e così via.
Immagina di partire da un singolo batterio che si riproduce ogni ora.

\begin{center}
  \begin{tikzpicture}[
      every node/.style={circle, draw=blu, fill=azzurro, inner sep=3pt, font=\small},
      level distance=1.1cm,
      level 1/.style={sibling distance=5cm},
      level 2/.style={sibling distance=2.5cm},
      level 3/.style={sibling distance=1.3cm},
      edge from parent/.style={draw=blu!60, thin}
    ]
    \node {1}
    child { node {2}
        child { node {4}
            child { node {8} }
            child { node {8} }
          }
        child { node {4}
            child { node {8} }
            child { node {8} }
          }
      }
    child { node {2}
        child { node {4}
            child { node {8} }
            child { node {8} }
          }
        child { node {4}
            child { node {8} }
            child { node {8} }
          }
      };
    % Etichette ore
    \node[draw=none, fill=none, right=0.3cm] at (4.5, 0)    {\small 0 ore};
    \node[draw=none, fill=none, right=0.3cm] at (4.5, -1.1) {\small 1 ora};
    \node[draw=none, fill=none, right=0.3cm] at (4.5, -2.2) {\small 2 ore};
    \node[draw=none, fill=none, right=0.3cm] at (4.5, -3.3) {\small 3 ore};
  \end{tikzpicture}
\end{center}

Quanti batteri ci sono dopo 5 ore? Bisogna moltiplicare 2 per sé stesso 5 volte:
\[
  2 \cdot 2 \cdot 2 \cdot 2 \cdot 2 = 32
\]

Che calcoli lunghi! Esiste però un modo per abbreviare una \textbf{moltiplicazione
  di fattori tutti uguali}: si chiama \textbf{elevazione a potenza}.

\begin{center}
  \begin{tabular}{@{} c c r @{}}
    \toprule
    \textbf{Ore trascorse} & \textbf{Moltiplicazione}                      & \textbf{Numero di batteri} \\
    \midrule
    0                      & $2^0$                                         & 1                          \\
    1                      & $2^1$                                         & 2                          \\
    2                      & $2 \times 2 = 2^2$                            & 4                          \\
    3                      & $2 \times 2 \times 2 = 2^3$                   & 8                          \\
    4                      & $2 \times 2 \times 2 \times 2 = 2^4$          & 16                         \\
    5                      & $2 \times 2 \times 2 \times 2 \times 2 = 2^5$ & 32                         \\
    12                     & $2^{12}$                                      & 4096                       \\
    \bottomrule
  \end{tabular}
\end{center}

\begin{sapevi}
  \textbf{Lo sapevi?}

  Lo yogurt che mangi a colazione esiste grazie ai batteri! Alcuni tipi di batteri
  «buoni», detti \textit{fermenti lattici}, trasformano il latte in yogurt attraverso
  un processo di fermentazione. Senza la riproduzione batterica — e senza le potenze
  per descriverla — non potremmo capire quanto rapidamente avviene questo processo.
\end{sapevi}

% ============================================================
\subsection{Perché servono le potenze? Tre storie matematiche}
% ============================================================

Le potenze non sono un'invenzione recente: i matematici le hanno usate (senza
saperlo) per migliaia di anni. Ecco tre esempi tratti dalla storia.

\subsubsection*{La scacchiera del bramino}

\begin{sapevi}
  \textbf{L'invenzione del gioco degli scacchi}

  Secondo un antico racconto, il re Iadava perse il suo unico figlio durante una guerra
  e cadde in una profonda disperazione. Un giorno, un umile bramino gli donò una tavola
  divisa in 64 caselle: era la scacchiera. Il re apprezzò talmente il gioco da chiedere
  al bramino come volesse essere ricompensato.

  Il bramino rispose: «Non voglio oro, terreni o palazzi. Mi darai un chicco di grano
  per la prima casella, due per la seconda, quattro per la terza, e così via,
  raddoppiando la quantità ad ogni casella fino all'ultima.»

  Il re rise: sembrava una richiesta da sciocco. Ma i matematici di corte calcolarono
  la verità. Il numero totale di chicchi richiesti è:
  \[
    1 + 2 + 2^2 + 2^3 + \cdots + 2^{63} = 2^{64} - 1 = 18\,446\,744\,073\,709\,551\,615
  \]
  Circa $1{,}8 \times 10^{19}$ chicchi — più di tutta la produzione mondiale di grano
  per migliaia di anni. Il bramino non era affatto uno sciocco!

  \textit{Adattato da: Malba Tahan, \textup{L'uomo che sapeva contare}, Salani.}
\end{sapevi}

\subsubsection*{I sette gatti di Ahmes}

Il seguente problema è contenuto nel \textbf{Papiro di Ahmes}, un documento egizio
del 1700 a.C. conservato al British Museum di Londra — uno dei testi matematici
più antichi del mondo.

\begin{esempio}
  \textit{In ogni casa ci sono sette gatti.
    Ogni gatto acchiappa sette topi.
    Ogni topo mangia sette spighe.
    Ogni spiga dà sette heqat di grano.
    Quante cose ci sono in tutto?}

  \textit{Soluzione:}
  \begin{align*}
    \text{case}   & = 7 = 7^1                               \\
    \text{gatti}  & = 7 \times 7 = 7^2 = 49                 \\
    \text{topi}   & = 7 \times 7 \times 7 = 7^3 = 343       \\
    \text{spighe} & = 7^4 = 2\,401                          \\
    \text{heqat}  & = 7^5 = 16\,807                         \\[4pt]
    \text{Totale} & = 7^1 + 7^2 + 7^3 + 7^4 + 7^5 = 19\,607
  \end{align*}
\end{esempio}

\subsubsection*{La filastrocca di Camogli}

Anche nella tradizione italiana compaiono moltiplicazioni ripetute:

\begin{center}
  \textit{Per una strada che mena a Camogli passava un uomo con sette mogli;}\\
  \textit{ed ogni moglie aveva sette sacche, ed in ogni sacca aveva sette gatte;}\\
  \textit{ed ogni gatta sette gattini.}\\[4pt]
  \textit{Fra gatti, gatte, sacche e mogli, in quanti andavano, dite, a Camogli?}
\end{center}

\begin{esempio}
  \textit{Soluzione:}
  \begin{align*}
    \text{mogli}   & = 7                               \\
    \text{sacche}  & = 7 \times 7 = 7^2 = 49           \\
    \text{gatte}   & = 7 \times 7 \times 7 = 7^3 = 343 \\
    \text{gattini} & = 7^4 = 2\,401                    \\[4pt]
    \text{Totale}  & = 7 + 7^2 + 7^3 + 7^4 = 2\,800
  \end{align*}
\end{esempio}

\begin{attenzione}
  Tutte e tre le situazioni — batteri, scacchiera, filastrocca — descrivono lo stesso
  fenomeno matematico: una quantità che si \textbf{moltiplica ripetutamente} per uno
  stesso numero. Le potenze sono lo strumento esatto per descrivere questo tipo di crescita.
\end{attenzione}

% ============================================================
\section{Che cos'è una potenza?}
% ============================================================

Quando si moltiplica un numero per sé stesso più volte, si può usare una scrittura
abbreviata chiamata \textbf{potenza}.

\begin{definizione}
  La \textbf{potenza} $a^n$ è il prodotto di $n$ fattori uguali al numero $a$:
  \[
    a^n = \underbrace{a \times a \times a \times \cdots \times a}_{n \text{ volte}}
  \]
  Il numero $a$ si chiama \textbf{base}, il numero $n$ si chiama \textbf{esponente}.
  Si legge: «$a$ elevato alla $n$» oppure «$a$ alla $n$».
\end{definizione}

\subsection{Le parti di una potenza}

\begin{esempio}
  Considera la potenza $3^4$:
  \[
    \underbrace{3}_{\text{base}}^{\!\!\!\underbrace{4}_{\text{esponente}}}
    = 3 \times 3 \times 3 \times 3 = 81
  \]
  Il risultato $81$ è detto \textbf{valore} della potenza.
\end{esempio}

\begin{esempio}
  Alcuni esempi di calcolo:
  \begin{align*}
    2^5  & = 2 \times 2 \times 2 \times 2 \times 2 = 32 \\
    5^3  & = 5 \times 5 \times 5 = 125                  \\
    10^4 & = 10 \times 10 \times 10 \times 10 = 10\,000
  \end{align*}
\end{esempio}

% ============================================================
\subsection{Casi particolari: esponente 0 e 1}
% ============================================================

\begin{regola}
  Per qualsiasi numero $a$ con $a \neq 0$ valgono le seguenti regole:
  \[
    a^1 = a \qquad \text{(un solo fattore uguale alla base)}
  \]
  \[
    a^0 = 1 \qquad \text{(qualsiasi numero diverso da zero elevato a zero vale 1)}
  \]
\end{regola}

\begin{esempio}
  \begin{align*}
    7^1   & = 7   \\
    7^0   & = 1   \\
    100^0 & = 1   \\
    100^1 & = 100
  \end{align*}
\end{esempio}

\begin{attenzione}
  \textbf{Caso speciale: $0^0$}

  Il valore di $0^0$ è una questione matematica delicata: in alcuni contesti
  si considera uguale a $1$, in altri è considerato indeterminato.
  A livello di scuola media, questa situazione non viene di solito incontrata,
  ma è importante sapere che $0^0$ non si comporta come le altre potenze di zero.

  Regola pratica: nelle espressioni che incontrerai, la base e l'esponente
  non saranno mai contemporaneamente zero.
\end{attenzione}

% ============================================================
\section{Le potenze nelle espressioni aritmetiche}
% ============================================================

\begin{regola}
  Nelle espressioni aritmetiche, l'\textbf{elevamento a potenza} ha
  \textbf{precedenza} rispetto a moltiplicazioni, divisioni, addizioni e sottrazioni.

  L'ordine completo delle operazioni è:
  \begin{enumerate}[leftmargin=2em]
    \item Operazioni tra parentesi (dall'interno verso l'esterno)
    \item \textbf{Elevamento a potenza}
    \item Moltiplicazioni e divisioni (da sinistra verso destra)
    \item Addizioni e sottrazioni (da sinistra verso destra)
  \end{enumerate}
\end{regola}

\begin{esempio}
  Calcola: $2 + 3 \times 2^3$

  \textit{Soluzione:}
  \begin{align*}
    2 + 3 \times 2^3
     & = 2 + 3 \times 8 &  & \text{(prima la potenza: } 2^3 = 8\text{)} \\
     & = 2 + 24         &  & \text{(poi la moltiplicazione)}            \\
     & = 26             &  & \text{(infine l'addizione)}
  \end{align*}
\end{esempio}

\begin{esempio}
  Calcola: $(2 + 3)^2 - 4 \times 2$

  \textit{Soluzione:}
  \begin{align*}
    (2 + 3)^2 - 4 \times 2
     & = 5^2 - 4 \times 2 &  & \text{(prima la parentesi)}     \\
     & = 25 - 4 \times 2  &  & \text{(poi la potenza)}         \\
     & = 25 - 8           &  & \text{(poi la moltiplicazione)} \\
     & = 17               &  & \text{(infine la sottrazione)}
  \end{align*}
\end{esempio}

\begin{attenzione}
  \textbf{Errore comune: confondere $(-3)^2$ e $-3^2$}

  \begin{align*}
    (-3)^2 & = (-3) \times (-3) = 9 \qquad \text{(la base è $-3$)}                \\
    -3^2   & = -(3 \times 3) = -9 \qquad \text{(si eleva prima $3$, poi si nega)}
  \end{align*}

  Le parentesi cambiano completamente il risultato!
\end{attenzione}

% ============================================================
\section{Proprietà delle potenze}
% ============================================================

Le proprietà che seguono valgono per basi e esponenti appartenenti
ai numeri naturali (con le restrizioni indicate caso per caso).

% ============================================================
\subsection{Prodotto di potenze con la stessa base}
% ============================================================

\begin{regola}
  Il prodotto di due potenze con la \textbf{stessa base} si calcola mantenendo
  la base e \textbf{sommando gli esponenti}:
  \[
    a^b \times a^c = a^{b+c}
  \]
\end{regola}

\begin{esempio}
  \[
    3^2 \times 3^4 = 3^{2+4} = 3^6 = 729
  \]
  \textit{Verifica:} $\;3^2 \times 3^4 = 9 \times 81 = 729$ \quad\textit{(corretto)}
\end{esempio}

\begin{attenzione}
  Questa proprietà vale \textbf{solo} se le due potenze hanno la stessa base.
  Non è possibile semplificare $2^3 \times 5^3$ usando questa regola,
  perché le basi sono diverse.
\end{attenzione}

% ============================================================
\subsection{Quoziente di potenze con la stessa base}
% ============================================================

\begin{regola}
  Il quoziente di due potenze con la \textbf{stessa base} si calcola mantenendo
  la base e \textbf{sottraendo l'esponente del divisore}:
  \[
    a^c \div a^b = a^{c-b} \qquad \text{con } a \neq 0 \text{ e } c \geq b
  \]
\end{regola}

\begin{esempio}
  \[
    5^7 \div 5^3 = 5^{7-3} = 5^4 = 625
  \]
  \textit{Verifica:} $\;5^7 \div 5^3 = 78125 \div 125 = 625$ \quad\textit{(corretto)}
\end{esempio}

% ============================================================
\subsection{Potenza di una potenza}
% ============================================================

\begin{regola}
  La \textbf{potenza di una potenza} si calcola mantenendo la base e
  \textbf{moltiplicando gli esponenti}:
  \[
    (a^b)^c = a^{b \times c}
  \]
\end{regola}

\begin{esempio}
  \[
    (2^3)^4 = 2^{3 \times 4} = 2^{12} = 4096
  \]
  \textit{Verifica:} $\;2^3 = 8$ e $\;8^4 = 4096$ \quad\textit{(corretto)}
\end{esempio}

\begin{attenzione}
  Non confondere $(a^b)^c$ con $a^{(b^c)}$:
  \[
    (2^3)^4 = 2^{12} = 4096 \qquad \text{ma} \qquad 2^{(3^4)} = 2^{81}
  \]
  Il secondo valore è enormemente più grande! Le parentesi sono decisive.
\end{attenzione}

% ============================================================
\subsection{Prodotto di potenze con lo stesso esponente}
% ============================================================

\begin{regola}
  Il prodotto di due potenze con il \textbf{stesso esponente} si calcola
  \textbf{moltiplicando le basi} e mantenendo l'esponente:
  \[
    b^a \times c^a = (b \times c)^a
  \]
\end{regola}

\begin{esempio}
  \[
    3^4 \times 2^4 = (3 \times 2)^4 = 6^4 = 1296
  \]
  \textit{Verifica:} $\;3^4 \times 2^4 = 81 \times 16 = 1296$ \quad\textit{(corretto)}
\end{esempio}

% ============================================================
\subsection{Quoziente di potenze con lo stesso esponente}
% ============================================================

\begin{regola}
  Il quoziente di due potenze con il \textbf{stesso esponente} si calcola
  \textbf{dividendo le basi} e mantenendo l'esponente:
  \[
    b^a \div c^a = \left(\frac{b}{c}\right)^a \qquad \text{con } c \neq 0
  \]
\end{regola}

\begin{esempio}
  \[
    6^3 \div 2^3 = \left(\frac{6}{2}\right)^3 = 3^3 = 27
  \]
  \textit{Verifica:} $\;6^3 \div 2^3 = 216 \div 8 = 27$ \quad\textit{(corretto)}
\end{esempio}

% ============================================================
\subsection{Riepilogo delle proprietà}
% ============================================================

\begin{center}
  \begin{tabular}{@{} l l @{}}
    \toprule
    \textbf{Proprietà}          & \textbf{Formula}                  \\
    \midrule
    Prodotto, stessa base       & $a^b \times a^c = a^{b+c}$        \\[4pt]
    Quoziente, stessa base      & $a^c \div a^b = a^{c-b}$          \\[4pt]
    Potenza di potenza          & $(a^b)^c = a^{b \times c}$        \\[4pt]
    Prodotto, stesso esponente  & $b^a \times c^a = (b \times c)^a$ \\[4pt]
    Quoziente, stesso esponente & $b^a \div c^a = (b \div c)^a$     \\
    \bottomrule
  \end{tabular}
\end{center}

% ============================================================
\section{Le potenze di 10 e i numeri grandi}
% ============================================================

\begin{definizione}
  Le \textbf{potenze di 10} sono potenze con base uguale a 10:
  \[
    10^1 = 10, \quad 10^2 = 100, \quad 10^3 = 1000, \quad \ldots
  \]
  Il numero di \textbf{zeri} nel risultato è uguale all'esponente.
\end{definizione}

\begin{center}
  \begin{tabular}{@{} c l r @{}}
    \toprule
    \textbf{Potenza} & \textbf{Forma estesa}   & \textbf{Nome} \\
    \midrule
    $10^0$           & $1$                     & uno           \\
    $10^1$           & $10$                    & dieci         \\
    $10^2$           & $100$                   & cento         \\
    $10^3$           & $1\,000$                & mille         \\
    $10^4$           & $10\,000$               & diecimila     \\
    $10^5$           & $100\,000$              & centomila     \\
    $10^6$           & $1\,000\,000$           & milione       \\
    $10^9$           & $1\,000\,000\,000$      & miliardo      \\
    $10^{12}$        & $1\,000\,000\,000\,000$ & bilione       \\
    \bottomrule
  \end{tabular}
\end{center}

\begin{esempio}
  Le potenze di 10 rendono facile moltiplicare numeri grandi:
  \[
    7 \times 10^4 = 7 \times 10\,000 = 70\,000
  \]
  \[
    3{,}5 \times 10^6 = 3{,}5 \times 1\,000\,000 = 3\,500\,000
  \]
\end{esempio}

% ============================================================
\section{La notazione scientifica}
% ============================================================

In scienza e tecnologia si incontrano spesso numeri molto grandi (o molto piccoli)
che sarebbero scomodi da scrivere per esteso. La \textbf{notazione scientifica}
permette di esprimerli in forma compatta.

\begin{definizione}
  Un numero è scritto in \textbf{notazione scientifica} quando è espresso nella forma:
  \[
    a \times 10^n
  \]
  dove $a$ è un numero con \textbf{una sola cifra prima della virgola}
  (cioè $1 \leq a < 10$) e $n$ è un numero naturale.
\end{definizione}

\begin{esempio}
  La distanza tra la Terra e il Sole è circa $150\,000\,000\,\text{km}$.

  \textit{Conversione in notazione scientifica:}
  \begin{align*}
    150\,000\,000 & = 1{,}5 \times 100\,000\,000 \\
                  & = 1{,}5 \times 10^8
  \end{align*}

  Si conta: la virgola deve essere spostata di \textbf{8 posizioni} verso sinistra.
  \[
    \underbrace{1\overbrace{50\,000\,000}^{8 \text{ cifre}}}
  \]
\end{esempio}

\begin{esempio}
  La massa della Terra è circa $4\,600\,000\,000\,000\,000\,000\,000\,\text{t}$.

  \textit{Conversione in notazione scientifica:}
  \begin{align*}
    4\,600\,000\,000\,000\,000\,000\,000 & = 4{,}6 \times 10^{21}
  \end{align*}

  La cifra significativa è $4{,}6$; il numero di posizioni spostate è $21$.
\end{esempio}

\begin{regola}
  \textbf{Come convertire un numero in notazione scientifica:}
  \begin{enumerate}[leftmargin=2em]
    \item Scrivi il numero con la virgola dopo la \textbf{prima cifra} non nulla.
    \item Conta di quante posizioni hai spostato la virgola: quel numero è l'esponente $n$.
    \item Scrivi il risultato nella forma $a \times 10^n$.
  \end{enumerate}
\end{regola}

\begin{sapevi}
  \textbf{Lo sapevi?}

  La notazione scientifica è usata ogni giorno da astronomi, fisici e ingegneri.
  La distanza dalla Terra alla stella più vicina (Proxima Centauri) è circa
  $4{,}0 \times 10^{13}\,\text{km}$ — un numero con 13 zeri che sarebbe
  difficilissimo da leggere nella forma estesa!

  Anche i computer usano una forma analoga internamente per rappresentare numeri
  molto grandi o molto piccoli: si chiama \textit{virgola mobile} (in inglese
  \textit{floating point}).
\end{sapevi}

% ============================================================
\section*{Riepilogo del capitolo}
% ============================================================

\begin{itemize}[leftmargin=1.5em]
  \item Una \textbf{potenza} $a^n$ è la moltiplicazione ripetuta della base $a$
        per $n$ volte. Il numero $n$ si chiama \textbf{esponente}.
  \item $a^0 = 1$ per ogni $a \neq 0$;\quad $a^1 = a$ per ogni $a$.
  \item Nelle espressioni aritmetiche, l'\textbf{elevamento a potenza}
        si esegue \textbf{prima} di moltiplicazioni, divisioni, addizioni
        e sottrazioni.
  \item Le proprietà principali delle potenze:
        \begin{itemize}
          \item Stesso base, prodotto: si sommano gli esponenti.
          \item Stessa base, quoziente: si sottraggono gli esponenti.
          \item Potenza di potenza: si moltiplicano gli esponenti.
          \item Stesso esponente, prodotto: si moltiplicano le basi.
          \item Stesso esponente, quoziente: si dividono le basi.
        \end{itemize}
  \item Le \textbf{potenze di 10} hanno tanti zeri quant'è l'esponente.
  \item La \textbf{notazione scientifica} $a \times 10^n$ (con $1 \leq a < 10$)
        permette di scrivere in modo compatto numeri molto grandi.
\end{itemize}

% ============================================================
%  (Sezione riservata per future soluzioni degli esercizi)
% ============================================================
% \section*{Soluzioni}